\icode{N17J-MP-SI-L2-261004-007}
\nceeID[6]{2017北京}{6}
2017年年初，我国研制的“大连光源”—极紫外自由电子激光装置，发出了波长在\qty{100}{nm}(\qty{1}{nm}=\,\qty{e-9}{m})附近连续可调的世界上首个最强的极紫外激光脉冲，大连光源因其光子的能量大、密度高，可在能源利用、光刻技术、雾霾治理等领域的研究中发挥重要作用。一个处于极紫外波段的光子所具有的能量可以电离一个分子，但又不会把分子打碎。据此判断，能够电离一个分子的能量约为（取普朗克常量\(h =\) \,\qty{6.63e-34}{J\cdot s}，真空光速\(c =\) \,\qty{3e8}{m/s}）\paren[B]
\begin{choices}
	\item \qty{e-21}{J}
	\item \qty{e-18}{J}
	\item \qty{e-15}{J}
	\item \qty{e-12}{J}
\end{choices}
\begin{answer}
	B\qquad 提示：\(E = h\nu = \frac{hc}{\lambda}\)
\end{answer}